\section{Introduction}
\label{sec:intro}

Large language models (LLMs) are compared by running them on static benchmarks and
reporting an average. This is expensive---agentic and code benchmarks execute hidden
tests or whole repositories for every item---and statistically wasteful, because items
differ enormously in how much they reveal about a system
\citep{lalor2016building,rodriguez2021evaluation}. Computerized adaptive testing (CAT),
the psychometric answer to the same problem for human examinees
\citep{lord1980applications,wainer2000cat}, has therefore been transferred to LLMs: item
response theory (IRT) models fitted on the item-level verdicts of previously evaluated
systems are used to pick anchor items \citep{vivek2024anchor,polo2024tinybenchmarks}, to
distill benchmarks \citep{kipnis2025metabench}, or to administer items adaptively by
maximum Fisher information \citep{hofmann2025fluid,li2026atlas,zheng2026medcat,zhou2026lost}.

These methods calibrate on hundreds to thousands of systems and then treat the item
parameters as known. Many practical settings have only a handful of calibration systems:
a new or private benchmark, a freshly generated item pool, a domain where only a few open
models can be run, or an agentic benchmark where every evaluation is costly. With four to
eight calibrators, item parameters---especially discriminations---are poorly identified,
and adaptive selection then \emph{capitalizes on calibration error}: it preferentially
picks the items whose informativeness is most overestimated
\citep{vanderlinden2000capitalization,patton2013influence}. Psychometrics knows remedies
for this: expected response functions that integrate the item-parameter posterior
\citep{mislevy1994erf} and fully Bayesian adaptive testing
\citep{vanderlinden2020fast,fink2025accounting}. What is not known is how much these
remedies matter in the small-panel regime of LLM evaluation, against which competitors,
and with what consequences for the uncertainty that an adaptive test reports.

This paper studies that question with a calibration-marginal adaptive test, EOST-M: items
enter the test through their expected response function under an empirical-Bayes item
prior, are selected by expected information gain (EIG), and are graded by an independent
oracle. Two studies with analysis plans frozen before execution---the second designed after
an external review, with one target per family cluster and the strongest competitors we
could identify---lead to three messages.

\begin{enumerate}\itemsep2pt
\item \textbf{Small-panel calibration uncertainty is the real problem.} With 4--8
calibrators, plug-in Fisher-information CAT treats item parameters that a small panel
cannot determine---poorly located difficulties and inflated discriminations---as known,
and both its predictions and its intervals suffer.
\item \textbf{Accounting for that uncertainty fixes a substantial part of it.} Integrating
the item-parameter posterior (EOST-M) and constraining the discriminations (a regularized
Rasch-type test) are two remedies for the same failure. \StudyTwoIntroMsgTwo
\item \textbf{Once uncertainty is handled, the choice of information criterion matters
much less than first hypothesized.} EIG, maximum and posterior-weighted Fisher information
and the score-targeted criterion performed almost identically under the
calibration-marginal model (a descriptive comparison, not an equivalence test), and the
original hypothesis that information beats predictive entropy on held-out prediction is
consistent in sign but small and not supported by the frozen rule; for the benchmark score,
entropy differs more, as the one-step score-targeted analysis predicts.
\end{enumerate}

\paragraph{Contributions.}
\begin{enumerate}\itemsep1pt
\item \textbf{Theory} (Section~\ref{sec:theory}). The calibration-marginal model is the exact
Bayesian model given the calibrators' abilities and the hyperprior (Lemma~\ref{lem:exact}),
which locates it with respect to fully Bayesian CAT; an item whose difficulty has a
Gaussian posterior with standard deviation $s$ can never carry more than $2/(\pi s^2)$ of
Fisher information, for the logistic link and any other, with a general bound for other
posteriors (Proposition~\ref{prop:cap}); entropy, Fisher
information and EIG are related exactly and agree under a common discrimination; and a
closed form for one-step score-targeted design.
\item \textbf{Evidence} (Sections~\ref{sec:study1}--\ref{sec:cost}). Study~1: eight real
benchmarks from five domains (\NSystemsTotal\ systems), panels of 4--64 calibrators,
re-analysed with family clusters as units. Study~2: independent targets (one per strict
family cluster), competitors including a regularized Rasch-type test, a posterior-draw marginal Bayesian
CAT, ATLAS-style 3PL CAT, tinyBenchmarks, MINCE and Random:mg, a leave-one-calibrator-out
recalibration of score intervals, and a temporal validation on later SWE-bench
submissions. A cost analysis gives break-even points.
\item \textbf{Reproducibility.} A Docker-only pipeline over pinned public sources; every
number in the paper is a macro generated from hashed artifacts and the qualitative
statements about results are re-checked on the artifacts by an automatic test; both analysis plans were
hashed and committed before execution (Section~\ref{sec:design}).
\end{enumerate}

\paragraph{Scope.} All evidence is a replay of recorded verdicts produced by deterministic
oracles or by recorded aggregate verdicts (Section~\ref{sec:design}); a protocol for live
deployment is released (Section~\ref{sec:live}) but could not be run in our execution
environment. The framework
also allows models to generate their own items (self-testing); this was explored only in a
historical pilot program (Section~\ref{sec:pilot}) and is not part of the evidence.

\paragraph{Main findings.} \IntroFindings
